% ============================================================
% DS MOD08 -- SUJET B : loi binomiale, loi de Poisson, loi normale
% BTS ATI / CRSA 1re année -- durée 1 h 30
% Pôle Formation UIMM - CVDL - S. JAUBERT
% Corrigé et barème : MOD08_DS_Binomiale_Poisson_Normale_ATI_CRSA_Sujet_B_Corrige_v1_0 (sujet jumeau : Sujet_A)
% ============================================================
\documentclass[a4paper,11pt]{article}
\usepackage[utf8]{inputenc}
\usepackage[T1]{fontenc}
% babel-french absent sur certains postes : chargement conditionnel
\IfFileExists{french.ldf}{\usepackage[french]{babel}}{}
\usepackage{helvet}
\usepackage{courier}
\renewcommand{\familydefault}{\sfdefault}
\usepackage[top=3cm, bottom=2.5cm, left=2cm, right=2cm, headheight=2cm]{geometry}
\usepackage{amsmath,amssymb,amsfonts}
\usepackage{graphicx}
\usepackage[dvipsnames,table]{xcolor}
\usepackage{tikz}
\usepackage[most]{tcolorbox}
\usepackage{fancyhdr}
\usepackage{enumitem}
\usepackage{array,needspace}
\usepackage{titlesec}
\usepackage{hyperref}

\definecolor{uimmblue}{HTML}{005C85}
\definecolor{uimmred}{HTML}{D31326}
\definecolor{uimmlightblue}{HTML}{E8F1F8}
\definecolor{uimmgray}{HTML}{333333}
\definecolor{uimmlightgray}{HTML}{F5F5F5}
\hypersetup{colorlinks=true, linkcolor=uimmblue, urlcolor=uimmblue}

\titleformat{\section}{\Large\bfseries\color{uimmblue}}{}{0pt}{}[\vspace{2pt}{\color{uimmblue}\titlerule[1.2pt]}]
\titleformat{\subsection}{\large\bfseries\color{uimmblue}}{}{0pt}{}
\setlist[enumerate,1]{label={\color{uimmblue}\textbf{\arabic*.}}, leftmargin=*}
\setlist[enumerate,2]{label=\alph*), leftmargin=*}
\setlist[itemize,1]{label={\color{uimmblue}\textbullet}}

\newtcolorbox{exo}[1]{colback=white, colframe=uimmblue, fonttitle=\bfseries\color{white},
  coltitle=white, title={#1}, boxrule=1pt, arc=2pt,
  left=8pt, right=8pt, top=5pt, bottom=6pt, before skip=10pt, after skip=6pt}
\newtcolorbox{consigne}{colback=uimmlightgray, colframe=uimmgray, boxrule=0.8pt, arc=3pt,
  left=8pt, right=8pt, top=6pt, bottom=6pt}

\newcommand{\pts}[1]{\hfill\mbox{\color{uimmred}\small\textbf{(#1)}}}
% Zone de réponse lignée : #1 = nombre de lignes
\newcommand{\zone}[1]{\par\smallskip\noindent\begin{tikzpicture}
  \foreach \i in {1,...,#1}{\draw[uimmgray!30] (0,-0.85*\i) -- (\linewidth-1pt,-0.85*\i);}
  \end{tikzpicture}\par\medskip}
\newcommand{\EnTete}[2]{%
\pagestyle{fancy}\fancyhf{}\renewcommand{\headrulewidth}{0pt}
\fancyhead[L]{\raisebox{-0.5\height}{\includegraphics[height=1.5cm]{logo_uimm_placeholder.jpg}}}
\fancyhead[C]{\begin{minipage}{7cm}\centering
  {\small\color{uimmblue}\textbf{Pôle Formation UIMM - CVDL}}\\
  {\footnotesize\color{uimmgray} BTS ATI / CRSA 1\textsuperscript{re} année\\ Mathématiques}\end{minipage}}
\fancyhead[R]{{\small\color{uimmred}\textbf{MOD08 -- #1}}}
\fancyfoot[C]{{\color{uimmblue}\rule{\textwidth}{1pt}}\\[2pt]
  {\footnotesize\color{uimmgray} Pôle Formation UIMM -- Centre-Val de Loire \hfill S. JAUBERT \hfill #2 -- \thepage}}}
\EnTete{DS Sujet B}{Sujet B -- DS 1 h 30}
\newtcolorbox{rappel}{colback=uimmlightblue, colframe=uimmblue, fonttitle=\bfseries\color{white},
  coltitle=white, title={Rappels de cours}, boxrule=1pt, arc=2pt,
  left=8pt, right=8pt, top=5pt, bottom=6pt, before skip=8pt, after skip=8pt}

\begin{document}

\begin{tcolorbox}[colback=uimmred, colframe=uimmred, sharp corners, boxrule=0pt,
  left=14pt, right=14pt, top=10pt, bottom=10pt]
{\color{white}\LARGE\bfseries Devoir surveillé -- MOD08 \hfill Sujet B}\\[4pt]
{\color{white}\large Loi binomiale, loi de Poisson, loi normale \hfill Durée~: 1 h 30}
\end{tcolorbox}

\vspace{4pt}
\noindent{\small\renewcommand{\arraystretch}{1.6}%
\begin{tabular}{@{}>{\bfseries\color{uimmblue}}l@{\hspace{6pt}}p{8.6cm}@{\hspace{10pt}}>{\bfseries\color{uimmblue}}l@{\hspace{6pt}}p{3.2cm}@{}}
Nom, prénom~: & \dotfill & Groupe~: & \dotfill \\
 & & Note~: & \dotfill\ / 20 \\
\end{tabular}}

\begin{consigne}
\textbf{Consignes.} Calculatrice autorisée en mode examen, aucun document. Répondre sur le sujet, dans les zones prévues.
\begin{itemize}[leftmargin=12pt, itemsep=1pt, topsep=2pt]
  \item Toute loi utilisée doit être \textbf{nommée avec ses paramètres}.
  \item Pour chaque probabilité demandée~: écrire l'évènement à l'aide de la variable aléatoire, \textbf{puis} le résultat. Un résultat seul ne rapporte pas tous les points.
  \item Probabilités arrondies à $10^{-4}$. Le barème est donné à titre indicatif~: partie 1 sur 8 points, partie 2 sur 8 points, partie 3 sur 4 points.
\end{itemize}
\end{consigne}

\begin{rappel}
\textbf{Loi binomiale $\mathcal{B}(n\,;p)$.} On répète $n$ fois, de façon identique et indépendante, une épreuve à deux issues (succès de probabilité $p$). $X$ compte le nombre de succès~; elle prend les valeurs $0, 1, \dots, n$.
\[
P(X = k) = \binom{n}{k}\,p^{k}\,(1-p)^{n-k},\quad \binom{n}{k} = \frac{n!}{k!\,(n-k)!},\quad E(X) = np,\quad \sigma(X) = \sqrt{np(1-p)}.
\]
\textbf{Loi de Poisson $\mathcal{P}(\lambda)$.} Nombre d'évènements rares (pannes, défauts, alarmes) sur une durée ou une zone donnée. Valeurs $0, 1, 2, \dots$
\[
P(X = k) = e^{-\lambda}\,\frac{\lambda^{k}}{k!},\qquad E(X) = \lambda,\qquad \sigma(X) = \sqrt{\lambda}.
\]
\textbf{Approximation de $\mathcal{B}(n\,;p)$ par $\mathcal{P}(\lambda)$, avec $\lambda = np$.} Elle est légitime lorsque les \textbf{trois} conditions sont réunies~: \ $n \geq 30$, \ $p \leq 0{,}1$, \ $np \leq 10$.

\smallskip
\textbf{Loi normale $\mathcal{N}(\mu\,;\sigma)$.} $\mu$ est la moyenne, $\sigma$ l'\textbf{écart-type}. Règle des 68 -- 95 -- 99,7~:\\
$P(\mu-\sigma\leq X\leq\mu+\sigma)\approx 0{,}68$, \ $P(\mu-2\sigma\leq X\leq\mu+2\sigma)\approx 0{,}95$, \ $P(\mu-3\sigma\leq X\leq\mu+3\sigma)\approx 0{,}997$.

\smallskip
\textbf{Calculatrice.} «~Fdp~» (ou «~pd~») donne $P(X = k)$~; «~FRép~» (ou «~cd~») donne $P(X \leq k)$.\\
{\small TI~: \texttt{binomFdp(n,p,k)}, \texttt{binomFRép(n,p,k)}, \texttt{poissonFdp($\lambda$,k)}, \texttt{poissonFRép($\lambda$,k)}, \texttt{normalFRép(a,b,$\mu$,$\sigma$)}.\\
Casio~: \texttt{Bpd}, \texttt{Bcd}, \texttt{Ppd}, \texttt{Pcd}, \texttt{Ncd}. \quad Borne infinie pour la loi normale~: \texttt{-1E99} ou \texttt{1E99}.}
\end{rappel}

% ============================================================
\newpage
\section{Partie 1 -- Loi binomiale \hfill {\normalsize\color{uimmred}8 points}}
\begin{exo}{Remise en état de vérins pneumatiques}
L'atelier de maintenance d'un site de production remet en état les vérins pneumatiques des machines (changement des joints, contrôle de la tige). Après remise en état, chaque vérin passe un test d'étanchéité sur banc. L'historique montre que $8\,\%$ des vérins remis en état présentent une \textbf{fuite}. Le responsable de l'atelier prélève $20$ vérins parmi ceux remis en état dans l'année. Le nombre de vérins traités est assez important pour assimiler ce prélèvement à un tirage avec remise.\\
On note $X$ le nombre de vérins \textbf{qui fuient} dans le prélèvement.
\end{exo}
\begin{enumerate}
  \needspace{7.4cm}
  \item Justifier que $X$ suit une loi binomiale. Préciser ses paramètres. \pts{1,5 pt}
  \zone{6}
  \needspace{6.5cm}
  \item Calculer l'espérance $E(X)$ et l'écart-type $\sigma(X)$. Interpréter $E(X)$ dans le contexte. \pts{1,5 pt}
  \zone{5}
  \needspace{6.5cm}
  \item Calculer la probabilité que le prélèvement contienne exactement $2$ vérins qui fuient. Détailler le calcul à l'aide de la formule, \textbf{sans} utiliser les fonctions de loi de la calculatrice. \pts{1 pt}
  \zone{5}
  \needspace{4cm}
  \item Pour chacun des évènements suivants, traduire l'évènement à l'aide de $X$, puis calculer sa probabilité.
  \begin{enumerate}
    \needspace{4cm}
    \item Aucun vérin du prélèvement ne fuit. \pts{0,75 pt}
    \zone{3}
    \needspace{4.8cm}
    \item Le prélèvement contient au moins un vérin qui fuit. \pts{0,75 pt}
    \zone{4}
    \needspace{5.6cm}
    \item Au moins $18$ vérins du prélèvement sont étanches. \pts{1 pt}
    \zone{5}
    \needspace{4.8cm}
    \item Le prélèvement contient au plus $3$ vérins qui fuient. \pts{1 pt}
    \zone{4}
  \end{enumerate}
  \needspace{5.6cm}
  \item Règle de suivi~: une \textbf{analyse des causes} est lancée si le prélèvement contient plus de $3$ vérins qui fuient. Calculer la probabilité de lancer cette analyse alors que la proportion de vérins qui fuient est bien de $8\,\%$. \pts{0,5 pt}
  \zone{4}
\end{enumerate}

% ============================================================
\newpage
\section{Partie 2 -- Loi de Poisson \hfill {\normalsize\color{uimmred}8 points}}
\subsection*{A. Appels de l'astreinte maintenance}
\begin{exo}{Contexte}
Dans une usine d'emballage qui tourne en continu, un technicien de maintenance est d'astreinte la nuit. Il reçoit des \textbf{appels de dépannage} (convoyeur bloqué, défaut de variateur, capteur déréglé), qui surviennent de façon aléatoire et indépendante. La GMAO montre qu'il y a \textbf{en moyenne $1{,}8$ appel par nuit}. On admet que le nombre d'appels au cours d'une nuit suit une loi de Poisson.
\end{exo}
\begin{enumerate}
  \needspace{5.6cm}
  \item Définir une variable aléatoire $Y$ adaptée à la situation. Donner sa loi et justifier la valeur de son paramètre. \pts{1 pt}
  \zone{3}
  \needspace{4cm}
  \item Calculer l'écart-type de $Y$. \pts{0,5 pt}
  \zone{2}
  \needspace{6.5cm}
  \item Calculer la probabilité que le technicien ne reçoive aucun appel dans la nuit, puis la probabilité qu'il en reçoive exactement $2$. Détailler les calculs à l'aide de la formule. \pts{1 pt}
  \zone{4}
  \needspace{4.8cm}
  \item Calculer la probabilité que le technicien reçoive plus de $3$ appels dans la nuit. \pts{1 pt}
  \zone{4}
  \needspace{5.4cm}
  \item Le responsable maintenance veut fixer un seuil $k$~: au-delà de $k$ appels dans une nuit, un second technicien est rappelé en renfort. Déterminer la plus petite valeur de $k$ telle que $P(Y \leq k) \geq 0{,}95$. Justifier par deux probabilités. \pts{1,5 pt}
  \zone{4}
\end{enumerate}

\newpage
\subsection*{B. Approximation d'une loi binomiale par une loi de Poisson}
\begin{exo}{Contexte}
Un fournisseur livre des câbles de capteurs à connecteur M12 surmoulé par lots de $300$ câbles. La probabilité qu'un câble soit défectueux est $0{,}008$. Le lot est assimilé à un tirage avec remise. On note $N$ le nombre de câbles défectueux dans un lot.
\end{exo}
\begin{enumerate}[resume]
  \needspace{4cm}
  \item Donner la loi exacte de $N$ et ses paramètres. \pts{0,5 pt}
  \zone{3}
  \needspace{5.6cm}
  \item Justifier, en vérifiant chacune des conditions, que l'on peut approcher la loi de $N$ par une loi de Poisson. Préciser son paramètre. \pts{1 pt}
  \zone{6}
  \needspace{4.8cm}
  \item À l'aide de cette loi de Poisson, calculer la probabilité qu'un lot contienne au plus $1$ câble défectueux. \pts{1 pt}
  \zone{5}
  \needspace{4.8cm}
  \item Calculer la même probabilité avec la loi exacte. Comparer les deux résultats et conclure sur la qualité de l'approximation. \pts{0,5 pt}
  \zone{5}
\end{enumerate}

% ============================================================
\newpage
\section{Partie 3 -- Loi normale \hfill {\normalsize\color{uimmred}4 points}}
\begin{exo}{Tension de sortie d'alimentations 24 V}
Les armoires électriques d'une cellule automatisée sont équipées d'alimentations à découpage 24\,V continu. On note $U$ la tension de sortie, en volts, mesurée à pleine charge sur une alimentation prise au hasard. On admet que $U$ suit la loi normale $\mathcal{N}(24{,}2\,;0{,}15)$.
\end{exo}
\begin{enumerate}
  \needspace{4.8cm}
  \item Calculer la probabilité que la tension soit comprise entre $24{,}0$\,V et $24{,}5$\,V. \pts{1 pt}
  \zone{4}
  \needspace{4.8cm}
  \item Certains capteurs de la cellule se mettent en défaut sous $23{,}9$\,V. Calculer la probabilité que la tension soit inférieure à $23{,}9$\,V. \pts{1 pt}
  \zone{4}
  \needspace{5.6cm}
  \item À l'aide de la règle des 68 -- 95 -- 99,7, donner l'intervalle qui contient environ $95\,\%$ des tensions. Ce résultat est-il cohérent avec celui de la question 2~? Justifier. \pts{1 pt}
  \zone{5}
  \needspace{4.8cm}
  \item Le site compte $500$ alimentations de ce modèle. Estimer le nombre moyen d'alimentations dont la tension est inférieure à $23{,}9$\,V. \pts{1 pt}
  \zone{4}
\end{enumerate}

\end{document}
